\documentclass[11pt]{article}
\usepackage[margin=1in]{geometry}
\usepackage{amsmath,amssymb,amsthm}
\usepackage{enumitem}
\usepackage{graphicx}
\usepackage[hidelinks]{hyperref}

\newtheorem{theorem}{Theorem}
\newtheorem{lemma}{Lemma}
\newtheorem{remark}{Remark}

\newcommand{\WPC}{\operatorname{WPC}}
\newcommand{\Sol}{\operatorname{Sol}}

\title{A Full Answer to the Weakly Precompact Domination Question}
\author{Run \texttt{fa\_banach\_001}, agent \texttt{agent\_lane\_16}}
\date{2026-06-28}

\begin{document}
\maketitle

\section*{Source Problem}

The source paper is Bo Xiang, Jinxi Chen, and Lei Li,
\emph{Weak precompactness in Banach lattices}, arXiv:2207.06038.
In Remark 3.7, after proving the result for $\sigma$-Dedekind complete
domain lattices, the authors ask whether the following assertions are
equivalent for arbitrary Banach lattices $E$ and $F$:
\begin{enumerate}[label=(\arabic*)]
\item every positive operator $S:E\to F$ dominated by a weakly precompact
positive operator is weakly precompact;
\item either
\begin{enumerate}[label=(\alph*)]
\item the norm of $E'$ is order continuous, or
\item every order interval in $F$ is weakly precompact.
\end{enumerate}
\end{enumerate}
The source paper states that the answer is affirmative when $E$ is
$\sigma$-Dedekind complete and leaves the general case open.

\begin{center}
\includegraphics[width=0.92\linewidth]{figures/open_problem_crop.png}
\end{center}

\section*{Proof Intuition}

The only missing direction in the source paper is the sufficiency of
condition (a) without an extra completeness or weak-order-unit hypothesis on
$E$. Weak precompactness is a sequential property: to show that $S(B_E)$ is
weakly precompact, it is enough to start with an arbitrary sequence
$(x_n)\subset B_E$ and find a weakly Cauchy subsequence of $(Sx_n)$.

The sequence $(x_n)$ lives in a much smaller lattice than $E$: the closed
sublattice $G$ generated by the countable set $\{x_n:n\in\mathbb N\}$.
This $G$ is separable. The hypothesis that $E'$ has order-continuous norm
passes to closed sublattices, so $G'$ has order-continuous norm. Separability
of $G$ gives a strictly positive functional on $G$, i.e. a weak order unit in
$G'$. Thus the source paper's Theorem 3.4 applies to the restricted
operators $S|_G\leq T|_G$, and it says that $S|_G$ is weakly precompact.
Applying this to the original sequence gives the desired weakly Cauchy
subsequence.

\section*{Auxiliary Lemmas}

\begin{lemma}[heredity to closed sublattices]
Let $E$ be a Banach lattice whose dual norm is order continuous. If $G$ is a
closed sublattice of $E$, then the norm of $G'$ is order continuous.
\end{lemma}

\begin{proof}
We use the standard criterion: the norm of $X'$ is order continuous if and
only if every norm-bounded disjoint sequence in $X$ is weakly null.

Let $(g_n)$ be a norm-bounded disjoint sequence in $G$. It is also a
norm-bounded disjoint sequence in $E$, hence it is weakly null in $E$ by the
criterion applied to $E$. If $\gamma\in G'$, extend $\gamma$ to a functional
$\widetilde\gamma\in E'$ by Hahn--Banach. Then
$\gamma(g_n)=\widetilde\gamma(g_n)\to0$. Thus $(g_n)$ is weakly null in $G$.
Applying the same criterion to $G$ shows that the norm of $G'$ is order
continuous.
\end{proof}

\begin{lemma}[separable lattices have dual weak units]
If $G$ is a separable Banach lattice, then $G'$ has a weak order unit.
\end{lemma}

\begin{proof}
If $G=\{0\}$ the assertion is trivial. Otherwise choose a norm-dense sequence
$(u_n)$ in the positive unit sphere of $G$. For each $n$, choose
$\varphi_n\in (G')_+$ with $\|\varphi_n\|=1$ and $\varphi_n(u_n)=1$.
Set
\[
  \varphi=\sum_{n=1}^{\infty} 2^{-n}\varphi_n\in (G')_+.
\]
We claim that $\varphi$ is strictly positive. Let $0<g\in G$ and normalize
$v=g/\|g\|$. Choose $u_n$ with $\|u_n-v\|<1/2$. Then
\[
  \varphi_n(v)\geq \varphi_n(u_n)-\|\varphi_n\|\|u_n-v\|>1/2,
\]
so $\varphi(g)>0$. A strictly positive functional on $G$ is precisely a weak
order unit of the dual Banach lattice $G'$.
\end{proof}

\section*{Main Theorem}

\begin{theorem}
Let $E$ and $F$ be Banach lattices. The following assertions are equivalent.
\begin{enumerate}[label=(\arabic*)]
\item For all positive operators $S,T:E\to F$ with $0\leq S\leq T$, if $T$
is weakly precompact, then $S$ is weakly precompact.
\item Either the norm of $E'$ is order continuous or every order interval in
$F$ is weakly precompact.
\end{enumerate}
\end{theorem}

\begin{proof}
The implication $(1)\Rightarrow(2)$ is Theorem 3.6 of Xiang--Chen--Li.

Assume (2). If every order interval in $F$ is weakly precompact, then
$(1)$ follows directly from Theorem 3.3 of Xiang--Chen--Li: the solid hull
of every weakly precompact set in $F$ is weakly precompact, and
$S(B_E)\subseteq \Sol(T(B_E))$.

It remains to handle the case where the norm of $E'$ is order continuous.
Let $0\leq S\leq T:E\to F$ and suppose that $T$ is weakly precompact. We prove
that $S$ is weakly precompact by the sequential definition. Let
$(x_n)$ be an arbitrary sequence in $B_E$, and let $G$ be the closed sublattice
of $E$ generated by $\{x_n:n\in\mathbb N\}$. Then $G$ is separable. By the
first lemma, the norm of $G'$ is order continuous; by the second lemma, $G'$
has a weak order unit.

The restricted operators $S_G=S|_G$ and $T_G=T|_G$ are positive operators
from $G$ into $F$, and $0\leq S_G\leq T_G$. Moreover $T_G$ is weakly
precompact, since
\[
  T_G(B_G)\subseteq T(B_E).
\]
The source paper's Theorem 3.4 applies to $G$ because $G'$ has
order-continuous norm and a weak order unit. Therefore $S_G$ is weakly
precompact. Hence the sequence $(Sx_n)=(S_Gx_n)$ has a weakly Cauchy
subsequence in $F$.

Since every sequence in $S(B_E)$ has a weakly Cauchy subsequence, $S(B_E)$ is
weakly precompact. Thus $S$ is weakly precompact, proving $(1)$.
\end{proof}

\section*{Verification Notes and Limitations}

The proof uses no computation. The two added lemmas are standard separable
Banach-lattice reductions; the rest of the argument invokes the source
paper's Theorems 3.3, 3.4, and 3.6. The result is a full answer to the
question as stated in Remark 3.7.

Bounded novelty check: the run's registry, solution, attempt, and proof-gap
indexes were searched for arXiv:2207.06038 and core phrases including
``weakly precompact'', ``dominated by a weakly precompact operator'', and
``order interval''. Web searches on 2026-06-28 for the exact title and exact
domination-problem phrases returned the source arXiv record and no later
independent resolution.

\section*{Human Review Recommendation}

The packet should be reviewed as a likely valid full solution. The main point
for review is the heredity lemma for order-continuity of the dual norm; once
that standard criterion is accepted, the proof is a direct separable
localization of the source paper's own theorem.

\begin{thebibliography}{9}

\bibitem{XiangChenLi}
Bo Xiang, Jinxi Chen, and Lei Li,
\emph{Weak precompactness in Banach lattices},
arXiv:2207.06038, 2022.

\bibitem{AB}
Charalambos D. Aliprantis and Owen Burkinshaw,
\emph{Positive Operators},
Springer, 2006.

\bibitem{MeyerNieberg}
Peter Meyer-Nieberg,
\emph{Banach Lattices},
Universitext, Springer, 1991.

\end{thebibliography}

\end{document}
